\section{Method}

\subsection{Agents and protocols}

Each agent is a single OpenAI chat completion (default
\texttt{gpt-4.1-mini}, with a model-cascade fallback to
\texttt{gpt-4o}, \texttt{gpt-4.1}, \texttt{gpt-4o-mini},
\texttt{gpt-4.1-nano}, and \texttt{gpt-3.5-turbo} when the per-day rate
limit on the primary is exhausted), seeded with a role
system message identifying its name (\textsc{orchestrator} or
\textsc{executor}\_$i$), the peer set, and the protocol-specific
scaffolding. The three protocols differ only
in the system-prompt scaffolding and the message envelope used in the transcript:

\begin{itemize}
\item \textbf{Native:} plain-text messages, no capability registry, agents addressed
      directly by name.
\item \textbf{MCP:} agents are told they communicate via JSON-RPC envelopes over a
      Model Context Protocol tools/list registry, and must briefly acknowledge the
      tool/role before producing content. Messages are wrapped in
      \texttt{<mcp:message from=... to=...>}.
\item \textbf{A2A:} agents have Agent Cards advertising their capabilities; replies are
      prefixed with \texttt{[from <sender> -> <recipient>]} so a router can dispatch.
\end{itemize}

\subsection{Topologies}

\begin{itemize}
\item \textbf{Centralized:} agent 0 (\textsc{orchestrator}) decomposes the task,
      dispatches one subtask per executor, then integrates their replies into a final
      answer.
\item \textbf{Chain:} a sequential pipeline agent\_0 $\rightarrow$ agent\_1
      $\rightarrow \cdots \rightarrow$ agent\_{N-1}, with each stage receiving only the
      previous stage's output.
\item \textbf{Fully connected:} a group chat where in each of two rounds every agent
      sees the full prior conversation and contributes; the system terminates when an
      agent emits \texttt{DONE: <answer>}.
\end{itemize}

\subsection{MAST annotation}

Each trial trace (the concatenated transcript) is annotated by
\texttt{agentdash.annotator} (overridden to model: \texttt{gpt-3.5-turbo}
because the upstream default of \texttt{o1-mini} is unavailable on most
OpenAI accounts), which returns the MAST failure-mode
dictionary\,\cite{mast2024}. We aggregate the 14 failure-mode codes into
the three coarse categories FC1 (\texttt{1.x}, specification issues), FC2 (\texttt{2.x},
inter-agent misalignment), and FC3 (\texttt{3.x}, task verification). Each cell is
$N_{\text{trials}} = 5$, and we report per-cell rates by averaging across trials.
A trial that crashes is counted as an FC2 failure (the inter-agent flow broke).

\subsection{Hyperparameters held fixed}

Worker model defaults to \texttt{gpt-4.1-mini};
\texttt{temperature}$=0.7$ for the headline sweep and $0.0$ for the
ablation; \texttt{max\_tokens}=400 per agent turn; \texttt{timeout}=60s
per OpenAI call; $N_{\text{trials}}=5$ per cell at the first pass with
follow-up replicates to 10--15 trials. Sweep~1 holds
$N_{\text{agents}}=3$ and \texttt{task\_type}=\textsc{code}.
